\documentclass{article}
\usepackage[T1]{fontenc}
\usepackage{lmodern}
\usepackage{graphicx}
\usepackage{amsmath,amssymb}
\usepackage[a4paper,margin=2cm]{geometry}
\usepackage[absolute,overlay]{textpos}
\setlength{\TPHorizModule}{1mm}
\setlength{\TPVertModule}{1mm}
\usepackage[indonesian]{babel}
\usepackage{hyperref}
\setlength{\emergencystretch}{2em}
% unit: Halaman 9

\begin{document}

% === Halaman 9 (llm) ===

Properti algoritma ElGamal:
\begin{enumerate}
  \item Bilangan prima, $p$ \hfill (tidak rahasia)
  \item Bilangan acak, $g$ ($g < p, g$ adalah akar primitif dari $p$) \hfill (tidak rahasia)
  \item Bilangan acak, $x$ ($2 \le x \le p - 2$) \hfill (rahasia, \textcolor{red}{kunci privat})
  \item $y = g^x \pmod{p}$ \hfill (tidak rahasia, \textcolor{red}{kunci publik})
  \item $m$ (plainteks) \hfill (rahasia)
  \item $a$ dan $b$ (cipherteks) \hfill (tidak rahasia)
\end{enumerate}

\end{document}
